\newcommand{\matiere}{Chimie : les alcanes}
% Préambule commun aux fiches d'exercices et de correction
\documentclass[11pt,a4paper]{article}
\usepackage[utf8]{inputenc}
\usepackage[T1]{fontenc}
\usepackage{lmodern}
\usepackage[french]{babel}
\usepackage[a4paper,margin=2.2cm,headheight=14pt]{geometry}
\usepackage{amsmath,amssymb}
\usepackage{icomma}
\usepackage[version=4]{mhchem}
\usepackage{chemfig}
\usepackage{tikz}
\usetikzlibrary{decorations.pathmorphing,arrows.meta,calc,patterns,angles,quotes}
\usepackage{enumitem}
\setlist[itemize]{label=\textbullet,nosep}
\setlist[enumerate]{nosep}
\usepackage{array}
\usepackage[most]{tcolorbox}
\usepackage{fancyhdr}
\usepackage{lastpage}
\usepackage{needspace}
\newcommand{\vv}[1]{\overrightarrow{#1}}

\setlength{\parindent}{0pt}
\setlength{\parskip}{5pt}
\setchemfig{atom sep=2.1em}

\pagestyle{fancy}
\fancyhf{}
\fancyhead[L]{\small Préparation du devoir}
\fancyhead[R]{\small\matiere}
\fancyfoot[C]{\small Page \thepage\ sur \pageref{LastPage}}
\renewcommand{\headrulewidth}{0.4pt}

% Titre d'un exercice
\newcommand{\exercice}[1]{%
  \par\bigskip\needspace{9\baselineskip}
  \begin{tcolorbox}[colback=gray!12,colframe=black,boxrule=0.5pt,arc=1pt,
    left=6pt,right=6pt,top=2pt,bottom=2pt]
  \textbf{#1}
  \end{tcolorbox}\par\smallskip}

% Titre d'une question dans la correction
\newcommand{\question}[1]{\par\medskip\textbf{#1}\par}

% Résultat encadré dans la correction
\newtcolorbox{resultat}{colback=gray!6,colframe=black,boxrule=0.4pt,arc=1pt,
  left=6pt,right=6pt,top=2pt,bottom=2pt}

% Données encadrées
\newtcolorbox{donnees}{colback=white,colframe=black!60,boxrule=0.4pt,arc=1pt,
  left=6pt,right=6pt,top=2pt,bottom=2pt}

% Page de titre
\newcommand{\entete}[2]{%
  \begin{center}
  {\LARGE\bfseries #1}\\[4pt]
  {\large #2}\\[2pt]
  {\normalsize Classe de Première}
  \end{center}}

\begin{document}

\entete{Série d'exercices}{Chimie : les alcanes}

\vspace{2pt}
Nom et prénom : \dotfill \quad Classe : \dotfill

\begin{donnees}
\textbf{Données valables pour toute la série.} Masses molaires atomiques en g/mol :
$M(\ce{C}) = 12$ ; $M(\ce{H}) = 1$ ; $M(\ce{O}) = 16$.
Volume molaire : $V_m = 22,4$ L/mol.
\end{donnees}

%-------------------------------------------------------------------------------
\exercice{Situation d'évaluation}
Tu découvres dans une revue scientifique que les alcanes brûlent dans l'air ou dans le dioxygène en
dégageant beaucoup de chaleur. Ils sont ainsi utilisés comme combustibles. La combustion complète de
$1,16$ g d'un alcane à chaîne ramifiée produit $3,52$ g de dioxyde de carbone et $1,8$ g d'eau.

Tu es sollicité pour déterminer le volume de dioxygène qu'a nécessité cette combustion complète.
\begin{enumerate}
  \item Donne la formule générale des alcanes.
  \item Écris l'équation-bilan générale de la combustion complète d'un alcane dans le dioxygène.
  \item Détermine la formule brute de l'alcane.
  \item Détermine le volume de dioxygène nécessaire à la combustion complète.
\end{enumerate}

%-------------------------------------------------------------------------------
\exercice{Exercice 1}
Écris les formules semi-développées des composés suivants :
\begin{enumerate}[label=\alph*.]
  \item 2-chloro-3-méthylbutane ;
  \item 2,3-dichloro-2-méthylpentane ;
  \item 2-bromo-4-chloro-3-éthyl-3-méthylhexane ;
  \item 1,1,2,2-tétrafluoroéthane.
\end{enumerate}

%-------------------------------------------------------------------------------
\exercice{Exercice 2}
Pour chacune des propositions suivantes, entoure la lettre \textbf{V} si la proposition est vraie ou
la lettre \textbf{F} si elle est fausse.

\renewcommand{\arraystretch}{1.35}
\begin{tabular}{|p{0.80\textwidth}|>{\centering\arraybackslash}p{0.04\textwidth}|>{\centering\arraybackslash}p{0.04\textwidth}|}
\hline
1) Un alcane est un composé organique renfermant uniquement des atomes de carbone et des atomes
d'hydrogène. & V & F \\ \hline
2) Un alcane non cyclique est un hydrocarbure dont la molécule est formée uniquement d'atomes de
carbone (C) et d'atomes d'hydrogène (H) liés entre eux par des liaisons covalentes simples. & V & F \\ \hline
3) L'alcane est ramifié si sa chaîne carbonée est acyclique et si au moins un atome de carbone du
squelette est relié à trois autres atomes de carbone. & V & F \\ \hline
4) Les alcanes cycliques ou cyclanes sont des alcanes dont la chaîne carbonée est ramifiée. & V & F \\ \hline
5) La formule générale brute des alcanes est \ce{C_nH_{2n-2}}. & V & F \\ \hline
6) La molécule d'alcane est dite saturée car elle ne peut pas additionner une autre molécule. & V & F \\ \hline
\end{tabular}

%-------------------------------------------------------------------------------
\exercice{Exercice 3}
Ton professeur de physique-chimie fait réagir du propane avec du dichlore en présence de lumière.
Il obtient deux produits monochlorés $A$ et $B$ isomères.
\begin{enumerate}
  \item Écris l'équation-bilan de la réaction.
  \item Écris la formule semi-développée et le nom de chacun des composés $A$ et $B$ obtenus.
\end{enumerate}

%-------------------------------------------------------------------------------
\exercice{Exercice 4}
Au cours d'une séance de TP, un groupe d'élèves de Première désire identifier un alcane $A$.
Pour cela, le groupe réalise dans un premier temps la combustion complète de $3,6$ g de l'alcane $A$.
Il obtient $11$ g de dioxyde de carbone et $5,4$ g d'eau. Ensuite, il réalise la monochloration
de $A$ et obtient un seul produit.

Tu es le rapporteur du groupe.
\begin{enumerate}
  \item Écris l'équation-bilan de la combustion complète des alcanes.
  \item Déduis-en la formule brute de l'alcane $A$.
  \item Donne les formules semi-développées et les noms de ses isomères.
  \item
  \begin{enumerate}[label=4.\arabic*.]
    \item Détermine, parmi les formules semi-développées précédentes, celle qui correspond à $A$.
          Donne son nom.
    \item Écris l'équation-bilan de la monochloration de $A$.
  \end{enumerate}
\end{enumerate}

%-------------------------------------------------------------------------------
\exercice{Exercice 5}
En vue d'écrire et d'exploiter correctement les équations-bilans de combustion des alcanes, il est
demandé à un groupe d'élèves dont tu fais partie de déterminer la composition en volume d'un mélange
gazeux de deux alcanes. Vous faites brûler dans un récipient sec $12$ cm$^3$ d'un mélange de propane et
de butane avec $100$ cm$^3$ de dioxygène. La combustion est complète. Après retour aux conditions
initiales, l'eau s'étant condensée, il reste $42$ cm$^3$ de dioxyde de carbone et $31$ cm$^3$ de
dioxygène.

Tu es le rapporteur du groupe.
\begin{enumerate}
  \item Écris l'équation-bilan de la réaction de combustion :
  \begin{enumerate}[label=1.\arabic*.]
    \item du propane avec le dioxygène ;
    \item du butane avec le dioxygène.
  \end{enumerate}
  \item Détermine le volume initial :
  \begin{enumerate}[label=2.\arabic*.]
    \item du propane ;
    \item du butane.
  \end{enumerate}
  \item Calcule le pourcentage en volume du propane et du butane dans le mélange initial.
\end{enumerate}

\end{document}
